% Two-stream instability as a test of PIC numerics -- final report
% Builds on Overleaf as-is (pdfLaTeX + BibTeX), or locally with `latexmk main.tex`.
% Missing figures render as labelled placeholder boxes; see figures/README.md.
% Audience: an undergraduate taking a first course in plasma physics.
\documentclass[10pt,twocolumn,letterpaper]{article}

\usepackage[letterpaper,margin=0.75in,columnsep=0.25in]{geometry}
\usepackage[T1]{fontenc}
\usepackage[utf8]{inputenc}
\usepackage{lmodern}
\usepackage{microtype}
\usepackage{amsmath,amssymb}
\usepackage{bm}
\usepackage{graphicx}
\usepackage{xcolor}
\usepackage{booktabs}
\usepackage{siunitx}
\usepackage[font=small,labelfont=bf]{caption}
\usepackage[numbers,sort&compress]{natbib}
\usepackage[colorlinks=true,allcolors=blue!50!black]{hyperref}

% ---- Notation -------------------------------------------------------------
\newcommand{\wb}{\omega_b}            % beam plasma frequency
\newcommand{\wpe}{\omega_{p}}         % total plasma frequency, wpe^2 = 2 wb^2
\newcommand{\vz}{v_0}                 % beam drift speed
\newcommand{\ppc}{N_\mathrm{ppc}}     % particles per cell
\newcommand{\dd}{\mathrm{d}}

% ---- Figure placeholder: \figorbox{file}{width}{label shown while missing} --
\newcommand{\figorbox}[3]{%
  \IfFileExists{figures/#1}%
    {\includegraphics[width=#2]{figures/#1}}%
    {\fbox{\parbox[c][0.55#2][c]{\dimexpr#2-2\fboxsep-2\fboxrule\relax}{\centering\small\textit{Placeholder:}\\ #3\\ \texttt{figures/\detokenize{#1}}}}}}

% ---- Margin notes for drafting (set \drafttrue -> \draftfalse to hide) -----
\newif\ifdraft \drafttrue
\newcommand{\todo}[1]{\ifdraft{\color{red}\textbf{[TODO:} #1\textbf{]}}\fi}

\title{\vspace{-1.5em}The Electron Two-Stream Instability as a Test of\\
       Particle-in-Cell Numerics}
\author{Harry Wu}
\date{\today}

\begin{document}
\maketitle

\begin{abstract}
\noindent
When two beams of electrons pass through each other, a small ripple in their density can grow
exponentially into a large electric wave. This is the two-stream instability. For cold beams of
equal density, pen-and-paper theory predicts exactly how fast each wavelength grows: the
fastest growth rate is $\gamma_{\max}=\wb/2$, where $\wb$ is the plasma frequency of one beam.
We reproduce this instability on a computer with the particle-in-cell (PIC) code
\textsc{Entity}, measure the growth rate, and compare it with the theory. Because we know the
right answer in advance, the comparison tells us how much error the simulation method itself
introduces. We change three numerical settings (the number of simulated particles, the grid
spacing, and the particle ``shape'') and see how each one affects the accuracy of the growth
rate and the amount of artificial heating.
\todo{one sentence on the headline result once runs are in.}
\end{abstract}

% ===========================================================================
\section{Introduction}
\label{sec:intro}

Picture two streams of electrons moving through each other in opposite directions, with
positive ions in the background so that the plasma is neutral overall. If the electron density
is perfectly uniform, nothing happens. But suppose a small ripple appears, with slightly more
electrons in some places and slightly fewer in others. The ripple produces a weak electric
field, which speeds up some electrons and slows down others. The electrons then bunch up even
more where the density was already high. More bunching gives a stronger field, which gives more
bunching, and the ripple grows exponentially. Growth stops only when the wave becomes strong
enough to trap electrons in its potential wells~\citep{BohmGross1949,Chen2016}. Because a small
disturbance grows on its own instead of dying away, this is called an \emph{instability}.

The two-stream instability is one of the simplest instabilities in plasma physics. When the
beams are cold (all electrons in a beam have the same velocity) and have equal density, the
growth rate can be calculated exactly. That makes it a useful test problem for computer
simulations~\citep{BirdsallLangdon1991}. Real plasmas contain enormous numbers of particles, so
a simulation cannot follow every one of them. A particle-in-cell (PIC) code instead follows a
much smaller number of ``macroparticles'' and computes the electric field on a grid of points
in space. Both shortcuts introduce errors. Having too few particles makes the simulation
noisy, and the grid can artificially heat the plasma~\citep{Langdon1970,Okuda1972}. Since we
know the exact answer for the two-stream instability, we can measure how large these errors
are.

The report is organized as follows. Section~\ref{sec:theory} derives the growth rate from the
cold-fluid equations. Section~\ref{sec:pic} explains how a PIC code works.
Section~\ref{sec:setup} describes our simulations, and Secs.~\ref{sec:results}
and~\ref{sec:discussion} compare the results with the theory.

% ===========================================================================
\section{Linear theory}
\label{sec:theory}

\paragraph{Setup.}
We take two cold electron beams, each with density $n_b$, moving along $x$ with velocities
$+\vz$ and $-\vz$. The ions are treated as a fixed, uniform positive background that keeps the
plasma neutral. Each beam has its own plasma frequency,
\begin{equation}
  \wb^2 = \frac{n_b e^2}{\varepsilon_0 m_e},
\end{equation}
which is the frequency at which that beam's electrons would oscillate if they were displaced
on their own. Because the total electron density is $2n_b$, the plasma frequency of all
electrons together is $\wpe^2=2\wb^2$.

\paragraph{Cold-fluid equations.}
We treat each beam $j$ as a fluid with density $n_j$ and velocity $v_j$. Each fluid obeys the
continuity equation and Newton's second law, and the two are coupled through Gauss's law:
\begin{align}
  \frac{\partial n_j}{\partial t} + \frac{\partial (n_j v_j)}{\partial x} &= 0, \\
  m_e\left(\frac{\partial v_j}{\partial t} + v_j\frac{\partial v_j}{\partial x}\right) &= -eE, \\
  \varepsilon_0\frac{\partial E}{\partial x} &= -e\,(n_1 + n_2 - 2n_b).
\end{align}

\paragraph{Linearizing.}
We write each quantity as its uniform value plus a small wave-like ripple, for example
$n_j = n_b + \tilde n_j\, e^{i(kx-\omega t)}$, and keep only terms that are first order in the
ripples. A beam moving at velocity $v_{0j}$ ($=\pm\vz$) sees the wave at the Doppler-shifted
frequency $\omega - k v_{0j}$. Newton's law gives the velocity ripple, and the continuity
equation turns it into a density ripple:
\begin{equation}
  \tilde v_j = \frac{-ie\tilde E}{m_e(\omega - k v_{0j})}, \qquad
  \tilde n_j = \frac{k n_b\,\tilde v_j}{\omega - k v_{0j}}.
\end{equation}
Substituting both density ripples into Gauss's law and dividing out $ik\tilde E$ gives the
\emph{dispersion relation}, which says which combinations of $\omega$ and $k$ are allowed:
\begin{equation}
  1 = \frac{\wb^2}{(\omega - k\vz)^2} + \frac{\wb^2}{(\omega + k\vz)^2}.
  \label{eq:disp}
\end{equation}

\paragraph{Solving for the growth rate.}
Multiplying through by both denominators turns Eq.~\eqref{eq:disp} into a quadratic equation for
$\omega^2$:
\begin{equation}
  \omega^2 = k^2\vz^2 + \wb^2 \pm \wb\sqrt{\wb^2 + 4k^2\vz^2}.
\end{equation}
Take the minus sign. When $k\vz<\sqrt2\,\wb$ this root is negative, so $\omega$ is purely
imaginary, $\omega = i\gamma$. The ripple then grows as $e^{i(kx-\omega t)}\propto e^{\gamma t}$,
with growth rate
\begin{equation}
  \gamma^2(k) = \wb\sqrt{\wb^2 + 4k^2\vz^2} - k^2\vz^2 - \wb^2 .
  \label{eq:gamma}
\end{equation}
Waves shorter than the cutoff (larger $k$) do not grow at all. Setting
$\dd\gamma^2/\dd(k^2)=0$ gives the fastest-growing wave:
\begin{equation}
  \gamma_{\max} = \tfrac{1}{2}\,\wb \quad\text{at}\quad k_{\max}\vz = \tfrac{\sqrt{3}}{2}\,\wb .
  \label{eq:gmax}
\end{equation}
In one growth time $1/\gamma$ the amplitude grows by a factor of $e$. This is called one
``$e$-folding.''

Equation~\eqref{eq:gamma} is the reference curve for every measurement below. It holds only
while the ripples are small, which is why it is called \emph{linear} theory. Once the wave
becomes large enough to trap electrons, exponential growth stops and the electrons swirl
around inside the wave's potential wells, forming vortices in a plot of position versus
velocity (phase space)~\citep{BirdsallLangdon1991}.

% ===========================================================================
\section{The particle-in-cell method}
\label{sec:pic}

\paragraph{Macroparticles.}
A real plasma has far too many electrons to simulate one by one. A PIC code instead uses
\emph{macroparticles}. Each macroparticle stands for a large number $w$ of real electrons, so
it carries $w$ times the charge and $w$ times the mass of one electron. The charge-to-mass
ratio $q/m$ is unchanged, so a macroparticle follows exactly the same path that a single
electron would. The electric field is stored only at a set of grid points spaced $\Delta x$
apart, and each grid cell holds $\ppc$ macroparticles on average.

Using few particles has a side effect: the density is noisy. If a cell holds about $\ppc$
particles, random fluctuations change that number by about $\sqrt{\ppc}$, so
$\delta n/n\sim\ppc^{-1/2}$. That is much noisier than a real plasma. For our problem the noise
is partly helpful, because it supplies the small initial ripple and we do not have to add one
by hand. It also has a cost. The instability starts from a larger amplitude, so it reaches
saturation sooner, and that leaves fewer $e$-foldings of clean exponential growth to
measure~\citep{BirdsallLangdon1991}.

\paragraph{The PIC cycle.}
The simulation advances in small time steps of length $\Delta t$. Each step repeats the same
four stages~\citep{BirdsallLangdon1991,HockneyEastwood1988}:
\begin{enumerate}\itemsep0pt
  \item \emph{Gather:} find the field at each particle from the nearby grid values.
  \item \emph{Push:} use the force to update each particle's velocity and position.
  \item \emph{Deposit:} add up the charge and current that the particles carry onto the grid.
  \item \emph{Solve:} use that current to update the fields on the grid.
\end{enumerate}

\paragraph{Particle shapes.}
Steps 1 and 3 need a rule that connects particles, which can sit anywhere, to grid points,
which are fixed. In PIC this rule is described by giving each macroparticle a small spatial
\emph{shape}. The simplest rule, \emph{nearest grid point}, assigns all of a particle's charge
to the closest grid point. This works but is noisy, because the charge jumps from one point to
the next as the particle moves. The next rule, \emph{cloud-in-cell}, treats the particle as a
small cloud one cell wide and splits its charge between the two nearest grid points in
proportion to how close it is to each. Higher-order shapes spread the charge smoothly over
three or four grid points. Smoother shapes average out more of the particle noise, especially
at very short wavelengths, but each particle then touches more grid points, so every step
takes longer. The same shape is used for gathering and depositing, which guarantees that a
particle does not push on itself.

\paragraph{Particle push.}
The push integrates Newton's second law, $m\,\dd\bm v/\dd t = q(\bm E+\bm v\times\bm B)$, using
the Boris method~\citep{Boris1970}, the standard choice in PIC codes. In our problem the
magnetic field is essentially zero and the particles move only along $x$, so the push reduces
to the familiar \emph{leapfrog} scheme:
\begin{align}
  v^{n+1/2} &= v^{n-1/2} + \frac{q}{m}E^{n}\,\Delta t, \\
  x^{n+1} &= x^{n} + v^{n+1/2}\,\Delta t .
\end{align}
Here a superscript $n$ means ``at time $n\Delta t$.'' Velocities are stored at times offset by half a
step from the positions, which makes the method more accurate than updating both at the same time.
Leapfrog has one important limit. A plasma naturally oscillates at $\wpe$, and if the time step
is too long to follow that oscillation the simulation blows up. The scheme is stable only for
$\wpe\Delta t<2$ and is accurate for $\wpe\Delta t\lesssim0.2$.
\todo{confirm which pusher twostream.toml selects (Boris vs.\ Vay) in Entity v1.3.3.}

\paragraph{Charge and current on the grid.}
After the push, each particle's motion during the step is recorded on the grid as an electric
current $J_x$. \textsc{Entity} does this in a way that conserves charge
exactly~\citep{Esirkepov2001,Umeda2003}: charge is never created or lost, so the continuity
equation $\partial_t\rho+\partial_x J_x=0$ holds on the grid up to computer round-off. As a
result, Gauss's law stays satisfied automatically. The code also smooths (\emph{filters})
the current over neighboring grid points a few times per step, which further suppresses
short-wavelength noise.
\todo{check Entity docs for the deposition scheme at each shape order.}

\paragraph{Updating the field.}
\textsc{Entity} solves the full Maxwell equations on a staggered grid (the Yee
mesh~\citep{Yee1966}). To stay stable, light must not cross more than one cell per step,
$c\Delta t<\Delta x$. In our one-dimensional problem with motion only along $x$ there is no
magnetic field, and Amp\`ere's law reduces to $\partial_t E_x=-J_x/\varepsilon_0$. The current
deposited by the particles therefore directly changes the electric field, and the problem
behaves like a purely electrostatic one.

\paragraph{Numerical heating.}
A PIC code does not conserve energy exactly, so the plasma can slowly heat up even when no
physical process heats it. The main cause is the grid. On a grid with spacing $\Delta x$, two waves whose wavenumbers differ by a multiple of
$2\pi/\Delta x$ have exactly the same values at every grid point, so the grid cannot tell them
apart. Short-wavelength particle noise can therefore masquerade as a longer wave (this is
called \emph{aliasing}). This matters
when the Debye length is much smaller than a cell, $\lambda_D\ll\Delta x$. In that case
aliasing feeds energy into the particles, and the plasma heats until $\lambda_D$ grows to about
$\Delta x$. This effect is known as the finite-grid instability~\citep{Langdon1970,Okuda1972}.
Because our beams are nearly cold, it is the main source of artificial heating we expect, and
it should be weaker for smoother particle shapes and finer grids. Heating also matters for the
physics. Warm beams have a spread of velocities $v_{th}$, which smears out the bunching and
lowers $\gamma$ below the cold value once $k v_{th}$ becomes comparable to $\gamma$.

% ===========================================================================
\section{Simulation setup}
\label{sec:setup}

\paragraph{Units.}
The code uses \emph{code units}: speeds are measured in units of the speed of light ($c=1$),
and lengths are measured in units chosen by the input file. A reference density $n_0$ defines a
reference plasma frequency $\omega_{p0}$ and skin depth $d_0=c/\omega_{p0}$. We set $d_0=10$, so
$\omega_{p0}=0.1$. What matters physically are dimensionless ratios like $\gamma/\wb$ and
$k\vz/\wb$. The table lists both the raw code values and these ratios.

\paragraph{Species.}
We use the \textsc{Entity} PIC code~\citep{Entity} (v1.3.3, \texttt{streaming} problem
generator, input file \texttt{twostream.toml}) in one spatial dimension. The boundaries are
periodic: a particle leaving the right edge re-enters on the left. The simulation contains four
particle species. Two are the electron beams. The other two are positron species that are
never moved, so they act as the fixed positive background that plays the role of the ions. The
beam speed is set by the \emph{four-velocity} $u_x=\gamma_L v_x/c=\pm0.1$, where $\gamma_L$ is
the relativistic Lorentz factor. At this low speed $\gamma_L\approx1$, so $\vz\approx0.1\,c$
and the beams are essentially non-relativistic.

\paragraph{Parameters.}
The input \texttt{densities = [0.5, 0.5]} sets the combined density of each electron--positron
pair, so each beam has $n_b=0.25\,n_0$. Then $\wb=\omega_{p0}\sqrt{0.25}=0.05$, and
Eq.~\eqref{eq:gmax} predicts $\gamma_{\max}=0.025$ at $k_{\max}\approx0.435$, which is a
wavelength of $\lambda_{\max}=2\pi/k_{\max}\approx14.4$.

Because the box is periodic, only waves that fit a whole number $m$ of wavelengths into the box
length $L=100$ can exist, with $k=2\pi m/L$. About 7 wavelengths of the fastest-growing wave
fit in the box, and the $m=7$ mode has $k\vz/\wb=0.875$, very close to the theoretical peak at
$0.866$. The run lasts $\wb t=50$, which is about 25 $e$-foldings at $\gamma_{\max}$. That is
long enough to see the wave grow out of the noise and then saturate.

Table~\ref{tab:params} lists the baseline settings. In each scan we change one of $\ppc$,
$\Delta x$, or the shape and keep the others fixed. The baseline grid is very fine: the beam
Debye length is $\lambda_D=v_{th}/\wb=0.2$, and a cell is only $\Delta x/\lambda_D\approx0.04$
of that. To reach the regime $\lambda_D\ll\Delta x$ where finite-grid heating appears, the
resolution scan therefore has to make the cells more than ten times larger.
\todo{update if the production runs change any of these values.}

\begin{table}[t]
  \centering
  \caption{Baseline parameters from \texttt{twostream.toml}, in code units ($c=1$). The right
  column gives the same quantity as a dimensionless ratio.}
  \label{tab:params}
  \small
  \begin{tabular}{@{}lll@{}}
    \toprule
    Parameter & Value & Normalized \\
    \midrule
    Beam four-velocity $u_x$            & $\pm0.1$            & $\vz\approx0.0995\,c$ \\
    Beam density $n_b$                  & $0.25\,n_0$         & $\wb=0.05$ \\
    Temperature $T/m_ec^2$              & $10^{-4}$           & $v_{th}/\vz\approx0.1$ \\
    Box length $L$                      & 100                 & $L\wb/\vz\approx50$ \\
    Cells                               & 12288               & $\Delta x/\lambda_D\approx0.04$ \\
    Time step (CFL $=0.5$)              & $\Delta t\approx0.0041$ & $\wb\Delta t\approx2\times10^{-4}$ \\
    Particles per cell $\ppc$           & 16                  & \\
    Current filter passes               & 4                   & \\
    Run time                            & 1000                & $\wb t=50$ \\
    Shape order                         & \todo{}             & \\
    \bottomrule
  \end{tabular}
\end{table}

\paragraph{Measuring $\gamma(k)$.}
At each output time we take the Fourier transform of $E_x(x,t)$ in $x$. This splits the field
into the allowed box waves and gives the amplitude $|\tilde E_k(t)|$ of each one. While a wave
grows exponentially, $|\tilde E_k|\propto e^{\gamma t}$, so a plot of $\ln|\tilde E_k|$ versus
$t$ is a straight line with slope $\gamma(k)$. We fit a line over the time window after the wave
has risen above the noise and before it saturates. To estimate the uncertainty, we repeat the
fit for several reasonable choices of window and take the spread of the results.
Figure~\ref{fig:growth} shows this for the fastest-growing wave.

\begin{figure}[t]
  \centering
  \figorbox{growth_mode.pdf}{\columnwidth}{$\ln|\tilde E_k(t)|$ for $k\approx k_{\max}$, with fit}
  \caption{Amplitude of the fastest-growing wave versus time. The dashed line is the
  straight-line fit during exponential growth; its slope is $\gamma$. \todo{caption details.}}
  \label{fig:growth}
\end{figure}

% ===========================================================================
\section{Results}
\label{sec:results}

\paragraph{Growth rate versus wavenumber.}
Figure~\ref{fig:gammak} compares the measured $\gamma(k)$ with the theory of
Eq.~\eqref{eq:gamma}.
\todo{describe agreement, value of $\gamma_{\max}/\wb$ and $k_{\max}\vz/\wb$, behavior near
the cutoff $k\vz=\sqrt2\,\wb$.}

\begin{figure}[t]
  \centering
  \figorbox{gamma_vs_k.pdf}{\columnwidth}{measured $\gamma(k)/\wb$ vs.\ $k\vz/\wb$ with theory}
  \caption{Measured growth rate (points) compared with the cold-beam theory,
  Eq.~\eqref{eq:gamma} (line). \todo{caption details.}}
  \label{fig:gammak}
\end{figure}

\paragraph{Effect of the numerical settings.}
Figure~\ref{fig:conv} shows how the error in $\gamma_{\max}$ and the drift in total energy
depend on $\ppc$, $\Delta x$, and particle shape. We expect two different kinds of behavior.
Fewer particles mean a louder initial noise level, which scales as $\ppc^{-1/2}$. That should
shorten the window of clean exponential growth without changing its slope. A grid that is too
coarse to resolve $\lambda_D$ should instead heat the beams, and warmer beams grow more slowly,
so in that case the measured $\gamma$ itself should drop.
\todo{replace expectations with what the runs show.}

\begin{figure}[t]
  \centering
  \figorbox{convergence.pdf}{\columnwidth}{error in $\gamma_{\max}$ and energy drift vs.\ $\ppc$, $\Delta x$, shape}
  \caption{How the measured $\gamma_{\max}$ and the drift in total energy change with the
  numerical settings. \todo{caption details.}}
  \label{fig:conv}
\end{figure}

% ===========================================================================
\section{Discussion and conclusions}
\label{sec:discussion}

\todo{Summarize: (i) how well PIC recovers $\gamma(k)$; (ii) which parameter matters most for
accuracy; (iii) the cost trade-off between $\ppc$ and shape order; (iv) limitations, e.g.\ a
finite thermal spread reduces $\gamma$ relative to the cold result, and only box modes are
sampled.}

\bibliographystyle{unsrtnat}
\bibliography{refs}

\end{document}
